\section{\texorpdfstring{$\R$}{R}的基数}

\begin{proposition}{}{}
	$\left|\left\{0,1\right\}^{\N}\right|=\left|\left[0,1\right)\right|=\left|\R\right|=\left|\mathcal{P}\left(\N\right)\right|$.
\end{proposition}

\begin{theorem}{}{}
	$\left|\R\right|=2^{\aleph_{0}}$.
\end{theorem}

\begin{corollary}{}{}
	$\left|\R\right|>\aleph_{0}$.
\end{corollary}

\begin{definition}{连续统的势}{}
	我们称与$\R$等势的集合的级数具有连续统的势,将其记作$\mathfrak{c}$.
\end{definition}

\begin{proposition}{}{}
	$\R$中的每个非退化区间,$\R$中任意可数子集的补集,$\R\setminus\Q$都和$\R$等势.
\end{proposition}